\chapter{Conclusioni}
\label{ch:conclusioni}

Questa tesi ha affrontato una domanda operativa precisa: se il premio per il
rischio di volatilità sulle opzioni Bitcoin sia catturabile in modo
sistematico, al netto di costi di transazione realistici, attraverso una
strategia replicabile di vendita di opzioni con copertura delta. La risposta
è affermativa, e si articola in sei conclusioni.

La prima è che il premio per il rischio di volatilità su Bitcoin è reale e
catturabile in modo sistematico. La strategia Top Vega produce uno Sharpe
ratio in-sample di 1,75, statisticamente significativo all'1\%
($p = 0{,}007$), ottenuto attraversando il bear market del 2022, e al netto
di costi di transazione modellati in modo conservativo. Il premio, del resto,
è strutturalmente ampio: la volatilità implicita eccede quella realizzata nel
93\% delle giornate, con uno spread medio di 16,2 punti percentuali.

La seconda è che la strategia più semplice batte quelle più sofisticate.
La selezione dell'opzione a vega massimo, priva di parametri da stimare,
sopravvive là dove falliscono lo scoring basato sul gap tra densità, la
selezione Premium/DTE, la modalità multi-posizione, i filtri di regime e lo
stop loss. La semplicità non è un punto di partenza, ma il residuo di un
vaglio severo.

La terza è che il percorso dal complesso al semplice è esso stesso un
risultato. Documentare le cinque ipotesi scartate, ciascuna con un criterio
di rigetto quantitativo, è un contributo metodologico contro il
\emph{publication bias} che affligge la letteratura sulle strategie di
trading. In particolare, l'individuazione dell'artefatto da opzioni
illiquide con prezzo di uscita nullo, responsabile del 116\% dei profitti
apparenti di una delle strategie, è un monito di portata generale per
chiunque backtesti strategie in opzioni.

La quarta è che l'audit ha trasformato risultati apparentemente ottimi in
risultati credibili. Le distorsioni corrette nel corso del lavoro (uscite
valorizzate al bid anziché all'ask, costi assenti, opzioni illiquide con
prezzi fittizi) avrebbero prodotto una strategia illusoriamente redditizia.
Solo la loro rimozione ha reso i risultati finali difendibili.

La quinta è che il \emph{timing} del premio resta un problema aperto. Cinque
indicatori sono stati testati come possibili filtri, e nessuno ha mostrato
potere predittivo sul profitto dei singoli trade: il fattore determinante non
sono le condizioni al momento dell'ingresso, ma ciò che accade durante il
trade. Individuare un segnale di timing affidabile rimane una sfida per la
ricerca futura.

La sesta è che la validazione cross-asset conferma la natura strutturale del
meccanismo. Il win rate della sola gamba opzionaria è identico su Bitcoin e
su SPY (82\%), su un campione azionario quattro volte più ampio e su
quattordici anni che comprendono crisi ed episodi di stress. L'alpha è
proporzionale al livello della volatilità implicita: Bitcoin, con una
volatilità implicita quattro volte superiore a quella dell'indice azionario,
è il mercato in cui il premio è abbastanza ampio da superare i costi di
copertura. Il premio di volatilità è un fenomeno generale; Bitcoin è, allo
stato attuale, il luogo migliore per raccoglierlo.

Resta un'avvertenza prospettica. Il mercato delle opzioni su criptovalute è
giovane e in rapida maturazione; l'ingresso di capitale dedicato e di market
maker istituzionali potrebbe, nel tempo, comprimere il premio verso i livelli
dei mercati tradizionali. La finestra di opportunità documentata in questo
lavoro non è necessariamente permanente. La replicabilità operativa della
strategia, dimostrata dal sistema di paper trading, ne consente tuttavia il
monitoraggio continuo, così da riconoscere per tempo l'eventuale erosione del
premio che oggi rende Bitcoin un caso di studio unico.
